\title{QLoRA: Efficient Finetuning of Quantized LLMs}

\begin{abstract}
We introduce \textsc{QLoRA}, a memory-efficient finetuning strategy that enables the adaptation of large language models with up to 65B parameters on a single 48GB GPU, while maintaining the finetuning performance of standard 16-bit approaches. In \textsc{QLoRA}, gradients are propagated through a fixed, 4-bit quantized pretrained language model to Low Rank Adapters~(LoRA). The leading models developed with this method, termed \textbf{Guanaco}, surpass all prior open-source models on the Vicuna benchmark and achieve 99.3\% of ChatGPT's performance after only 24 hours of finetuning on a single GPU. The key techniques introduced by \textsc{QLoRA} to achieve high efficiency and maintain quality include: (a) the 4-bit NormalFloat (NF4), an information-theoretically optimal datatype for weights with a normal distribution, (b) Double Quantization, which further minimizes memory usage by quantizing the quantization constants themselves, and (c) Paged Optimizers, which address transient increases in memory demand. Leveraging \textsc{QLoRA}, we finetuned over 1,000 models and conducted comprehensive assessments of instruction-following and chatbot abilities across eight instruction datasets, a variety of model architectures (such as LLaMA and T5), and parameter sizes that exceed the practical limits of conventional finetuning methods (including 33B and 65B models). Our findings indicate that employing QLoRA for finetuning on concise, high-quality datasets enables state-of-the-art results, even for models smaller than previous leading systems. We offer a thorough evaluation of chatbot capabilities using both human assessments and GPT-4-based judgments, revealing that GPT-4 scoring offers a cost-effective and practical substitute for human evaluation. Additionally, we demonstrate that the reliability of existing chatbot benchmarks is questionable when it comes to evaluating chatbot performance. A curated failure analysis further illustrates the cases in which \textbf{Guanaco} underperforms relative to ChatGPT. All models, code, and CUDA kernels supporting 4-bit training are publicly available.\footnote{\url{https://github.com/artidoro/qlora} and \url{https://github.com/TimDettmers/bitsandbytes}}
\end{abstract}

\section{Introduction}

Adapting large language models (LLMs) through finetuning has proven to be a powerful strategy for enhancing their capabilities \citep{min2021metaicl, wei2021finetuned, ouyang2022training, wang2022super, wang2022self, liu2022few} and for instilling or suppressing specific behaviors as desired \citep{ouyang2022training,askell2021general,bai2022training}. Despite its efficacy, the computational demands of finetuning massive models remain a major bottleneck; standard 16-bit finetuning for a LLaMA model with 65 billion parameters~\citep{touvron2023llama} requires in excess of 780 GB of GPU memory. While advances in quantization techniques offer notable reductions in memory usage for inference \citep{dettmers2022llm, dettmers2022case, frantar2022gptq, xiao2022smoothquant}, these approaches typically fail when applied during training \citep{wortsman2023stable}. 

Here, we show that, contrary to previous assumptions, finetuning a quantized 4-bit model can be achieved with no loss in model performance. Our proposed approach, \textsc{QLoRA}, utilizes an innovative high-precision quantization procedure to convert a pretrained model to 4 bits, after which a compact set of Low-rank Adapter weights \citep{hu2021lora} is introduced. 

These adapter weights are optimized by propagating gradients through the quantized model parameters. 

\textsc{QLoRA} lowers the average GPU memory necessary to finetune a 65B parameter model from over 780GB to under 48GB, without sacrificing runtime speed or predictive accuracy when compared to traditional 16-bit finetuning. This development substantially broadens the practicality of LLM finetuning, making it feasible to train even the largest publicly available models on a single GPU. Leveraging \textsc{QLoRA}, we introduce the \textbf{Guanaco} model series: the second-best in this family achieves 97.8\% of ChatGPT's performance on the Vicuna~\citep{vicuna2023} benchmark, and can be trained within 12 hours using a consumer-grade GPU; with a single high-end GPU over 24 hours, our largest Guanaco model attains 99.3\%, virtually matching ChatGPT on Vicuna. When in operation, the most compact Guanaco configuration (7B parameters) consumes only 5 GB of memory and surpasses the 26 GB Alpaca model by more than 20 points on the Vicuna benchmark (see Table~\ref{tbl:vicuna_eval}).

\textsc{QLoRA} incorporates several novel techniques aimed at minimizing memory consumption while maintaining high performance: (1) \textbf{4-bit NormalFloat}, a quantization data type that is theoretically optimal for data with a normal distribution and achieves superior empirical performance compared to 4-bit Integer and 4-bit Float alternatives. (2) \textbf{Double Quantization}, an approach where the quantization coefficients themselves are quantized, resulting in an average memory saving of around 0.37 bits per parameter (which equates to roughly 3 GB for a model with 65B parameters). (3) \textbf{Paged Optimizers}, which leverage NVIDIA unified memory to circumvent the memory spikes associated with gradient checkpointing during the processing of mini-batches containing long sequences. These components are unified into an enhanced LoRA framework that incorporates adapters at every layer of the network, nearly eliminating the accuracy compromises observed in earlier studies.

Thanks to the resource efficiency of \textsc{QLoRA}, we are able to conduct a thorough investigation of instruction finetuning and chatbot effectiveness across model sizes that are not feasible with standard finetuning due to substantial memory demands. We thus train over 1,000 models utilizing various instruction tuning datasets, spanning model architectures and parameter counts from 80M to 65B. Beyond demonstrating that \textsc{QLoRA} matches the performance of 16-bit training (\S\ref{sec:comp_to_std}) and enables the development of the advanced chatbot \textbf{Guanaco} (\S\ref{sec:pushing}), our analysis also uncovers several important trends. Firstly, data quality proves to be significantly more impactful than the sheer size of the dataset; for example, a dataset containing 9,000 samples (OASST1) surpasses a much larger 450,000 sample dataset (FLAN v2, subsampled) in chatbot performance, even when both are tailored for general instruction following. Secondly, we demonstrate that excelling in the Massive Multitask Language Understanding (MMLU) benchmark does not necessarily translate to superior outcomes on the Vicuna chatbot benchmark, and vice versa—underscoring that for any specific task, dataset appropriateness outweighs its size.

In addition, we conduct a thorough evaluation of chatbot performance utilizing assessments from both human judges and GPT-4. Our methodology employs a tournament-style framework where different models compete directly by generating responses to the same prompt, and their outputs are then compared. The determination of the superior response in each pairwise contest is made either by GPT-4 or by human raters. Performance across all matchups is aggregated using Elo scores~\citep{elo1967proposed, elo1978rating}, providing an overall ranking of the chatbot systems. We observe a general consistency between GPT-4 and human-based rankings in the tournaments; however, we identify notable cases where their judgments diverge significantly. Thus, we emphasize that while model-driven evaluations offer a cost-effective alternative to human judgment, they are also accompanied by distinct limitations and uncertainties.

To complement our quantitative benchmarking, we provide a qualitative examination of the \textbf{Guanaco} models. This qualitative review sheds light on the strengths and weaknesses of the models, identifying behaviors not revealed by numerical benchmarks alone.

All model outputs, along with their corresponding human and GPT-4 annotations, are publicly available to enable further analysis by the community. We also release our full codebase and CUDA implementations, integrating our techniques within the Hugging Face transformers framework \citep{wolf2019huggingface} to ensure broad accessibility. Our open-source release comprises a suite of adapters for model sizes 7/13/33/65B, each trained on eight distinct instruction-following datasets, resulting in 32 unique, fine-tuned and publicly available models.

\section{Background}
\paragraph{Block-wise k-bit Quantization} Quantization refers to transforming input data from a high-information representation to one that encodes less information, typically by reducing numerical precision. This process often involves converting a data type with more bits, such as 32-bit floating point numbers, to one with fewer bits, such as 8-bit integers. To maximize the effective range of the target lower-bit datatype, the input values—usually organized in tensors—are rescaled to align with the new type’s range. This normalization often uses the absolute maximum value present in the original input. For example, converting an FP32 tensor into an Int8 tensor defined over $[-127, 127]$ is achieved as follows:

\begin{equation}
    \mathbf{X}^{ \text{Int8}} = \text{round}\left( \frac{127}{{\text{absmax}}(\mathbf{X}^{{\text{FP32}}})}\mathbf{X}^{{\text{FP32}}} \right) = \text{round}(c^{\text{FP32}}\cdot \mathbf{X}^{{\text{FP32}}}),
\end{equation}

where $c$ stands for the {\em quantization constant} or {\em quantization scale}. The original floating-point representation can be reconstructed via dequantization, defined as:

\begin{equation}
    \text{dequant}(c^{\text{FP32}}, \mathbf{X}^{\text{Int8}}) = \frac{\mathbf{X}^{\text{Int8}}}{c^{\text{FP32}}} = \mathbf{X}^{\text{FP32}}
\end{equation}

This technique faces challenges when the input tensor contains high-magnitude elements, or outliers, since these values skew the quantization range such that several quantization bins—i.e., specific bit assignments—are underutilized, with minimal or no values mapped to them. To address this issue with outliers, a prevalent strategy involves dividing the input tensor into smaller blocks, each of which undergoes separate quantization with an individual scaling constant $c$. Specifically, the input tensor $\mathbf{X}\in \mathbb{R}^{b\times h}$ is partitioned into $n$ sequential blocks, each of dimension $B$, by first flattening the tensor and then segmenting it linearly into ${n = ({b\times h} )/{B} }$ blocks. Each block is then quantized separately following Equation~1, resulting in a quantized tensor alongside $n$ unique scaling factors $c_i$.

\paragraph{Low-rank Adapters} Low-rank Adapter (LoRA) fine-tuning \citep{hu2021lora} introduces a strategy that minimizes memory consumption by introducing a compact collection of tunable parameters—adapters—while keeping the primary model parameters static. During stochastic gradient descent, gradients propagate through the unaltered pretrained weights to update only the adapter modules in pursuit of optimizing the loss objective. LoRA extends linear transformations by inserting an additional low-rank, factorized projection. For a standard linear mapping $\mathbf{X} \mathbf{W} = \mathbf{Y}$, with $\mathbf{X}\in  \mathbb{R}^{b\times h}$ and $\mathbf{W}\in  \mathbb{R}^{h\times o}$, LoRA modifies the output as:

\begin{equation}
    \mathbf{Y} = \mathbf{X}\mathbf{W} + s\mathbf{X}\mathbf{L}_1\mathbf{L}_2,
\end{equation}

where the matrices $\mathbf{L}_1\in  \mathbb{R}^{h\times r}$ and $\mathbf{L}_2\in  \mathbb{R}^{r\times o}$ represent the low-rank projections, and $s$ is a scaling factor.

\paragraph{Memory Requirement of Parameter-Efficient Finetuning} A key aspect to consider is the memory usage associated with LoRA during the training process, both regarding the quantity and scale of adapters implemented. Due to LoRA’s minimal memory consumption, it is feasible to incorporate a greater number of adapters in order to enhance performance, with only a marginal effect on overall memory usage. Although LoRA is categorized as a Parameter Efficient Finetuning (PEFT) technique, the principal contributor to memory consumption during LLM finetuning arises from activation gradients rather than the storage of LoRA’s learnable parameters. For instance, with a 7B LLaMA model finetuned on FLAN v2 at a batch size of 1, and LoRA weights comprising approximately 0.2\% of the original parameter count~\citep{hu2021lora,liu2022few}, the memory requirement for LoRA input gradients is 567 MB, whereas the LoRA parameters themselves only occupy 26 MB. Employing gradient checkpointing~\citep{chen2016training} reduces input gradient memory usage per sequence to roughly 18 MB, which is still more than the collective memory usage of all LoRA weights. By comparison, the 4-bit base model itself utilizes 5,048 MB. This comparison illustrates both the importance of gradient checkpointing and that aggressively minimizing LoRA parameter size yields limited memory savings. Consequently, adopting a greater number of adapters does not substantially elevate the total memory requirement for training (for comprehensive figures, see Appendix~\ref{app:memory}). As will be detailed further, this flexibility is vital for matching the performance achieved at full 16-bit precision.

\section{\textsc{QLoRA} Finetuning}

\textsc{QLoRA} enables precise 4-bit finetuning through two main innovations: 4-bit NormalFloat (NF4) quantization and Double Quantization. Furthermore, the introduction of Paged Optimizers addresses the challenge of memory spikes associated with gradient checkpointing, thereby mitigating the risk of out-of-memory issues that have previously hindered large model finetuning on a single device.

In the \textsc{QLoRA} framework, a low-precision storage format—typically 4-bit—is paired with a computation format, usually BFloat16. In practical terms, this setup entails that, whenever a \textsc{QLoRA} weight tensor is accessed, it is dequantized into BFloat16 and subsequently subjected to matrix multiplication at 16-bit precision.

The subsequent discussion covers the fundamental elements of \textsc{QLoRA} and concludes with a precise characterization of the approach.

\paragraph{4-bit NormalFloat Quantization} The NormalFloat (NF) type extends upon Quantile Quantization \citep{dettmers2022optimizers}, a theoretically optimal format that allocates an equal quantity of tensor values to each quantization interval. This is achieved by employing the empirical cumulative distribution function to determine quantiles within the input tensor.

However, a significant drawback of quantile quantization lies in the computational overhead of estimating quantiles. For this reason, accelerated quantile estimation methods, like SRAM quantiles~\citep{dettmers2022optimizers}, are commonly employed. These approximations, though efficient, tend to introduce substantial quantization error for outlier values, which can often be the most critical entries.

Nevertheless, if the input tensors originate from a distribution that is fixed modulo a quantization constant, one can circumvent the need for costly and inexact quantile estimates. In such situations, the uniformity of quantiles across tensors allows for efficient, exact computation of the quantization process.

Because the pretrained weights of neural networks tend to follow a zero-mean normal distribution with standard deviation $\sigma$ (refer to Appendix~\ref{app:normality}), it is possible to remap all weights to a standardized distribution by adjusting $\sigma$ such that the full distribution aligns precisely with the dynamic range of the intended data type. In our approach, we designate the data type’s permissible range as $[-1, 1]$. Consequently, we normalize both the quantiles of the data type and those of the neural network weights to this interval.

To derive the information-theoretically optimal data type for representing normal distributions (centered at zero, arbitrary standard deviation $\sigma$) within $[-1, 1]$, we carry out the following procedure: (1) identify the $2^k+1$ quantiles from the canonical $N(0, 1)$ distribution to produce a $k$-bit quantile-based quantization for normal distributions, (2) transform the resulting data type so that its values fall within $[-1, 1]$, (3) for a given input weight tensor, normalize its entries into $[-1, 1]$ by scaling with respect to its absolute maximum value.

After ensuring the value ranges of the weight tensor and data type coincide, standard quantization can be performed. The third step, in effect, adjusts the weight tensor’s standard deviation so it equals that of the $k$-bit data type. Specifically, the $2^k$ representative quantization levels $q_i$ for the data type are determined as:

\begin{equation}
q_i  = \frac{1}{2}\left( Q_X\left(\frac{i}{2^k+1}\right) + Q_X\left(\frac{i+1}{2^k+1}\right)\right),
\end{equation}

where $Q_X(\cdot)$ denotes the quantile function associated with the standard normal distribution $N(0,1)$. One issue with employing a symmetric k-bit quantization scheme is its inability to precisely represent zero, a key property necessary to quantize padding and other elements with value zero without introducing error. To guarantee that the discrete zeropoint is set to $0$ and to utilize the complete $2^k$ value range for a k-bit type, we construct an asymmetric data type by separately estimating the quantiles $q_i$ for the negative ($2^{k-1}$ values) and positive ($2^{k-1} + 1$ values) ranges. After forming these two sets of $q_i$, we merge them, ensuring to eliminate the duplicated zero present in both partitions. We refer to this data type as \emph{k-bit NormalFloat} (NFk), characterized by each quantization interval containing, in expectation, the same number of samples—rendering it information-theoretically optimal for data that follows a zero-centered normal distribution. The precise quantization values of NFk are detailed in Appendix~\ref{app:nf4}.

\paragraph{Double Quantization{}} 
We propose \emph{Double Quantization} (DQ), which entails quantizing the quantization parameters themselves, thereby yielding further memory efficiency. While the use of a small blocksize is essential for accurate 4-bit quantization as discussed by~\citet{dettmers2022case}, this practice comes with notable memory costs. For instance, with 32-bit quantization constants and a blocksize of 64 when quantizing $\mathbf{W}$, the associated quantization parameters contribute, on average, $32/64 = 0.5$ bits per parameter. The DQ method directly tackles the additional memory required for storing these constants.

In more detail, Double Quantization treats the first-stage quantization constants $c_2^{\text{FP32}}$ as inputs for a subsequent quantization procedure. This step produces quantized versions $c_2^{\text{FP8}}$ as well as a second layer of quantization constants $c_1^{\text{FP32}}$. For this secondary quantization, we utilize 8-bit floating-point numbers with a blocksize of 256, motivated by evidence from \citet{dettmers2022case} that shows no decline in performance at this precision. Because $c_2^{\text{FP32}}$ are strictly positive, we mean-center these values prior to quantization to enable the use of symmetric quantization schemes. Overall, assuming a blocksize of 64, the average memory required per parameter is reduced from $32/64=0.5$ bits down to $8/64 + 32/(64\cdot256) = 0.127$ bits, representing a memory saving of 0.373 bits per parameter.

\paragraph{Paged Optimizers} leverage NVIDIA’s unified memory~\footnote{\tiny \url{https://docs.nvidia.com/cuda/cuda-c-programming-guide}} capability, which transparently handles the transfer of memory pages between CPU and GPU. This mechanism ensures seamless GPU computation, even in circumstances where the GPU experiences intermittent out-of-memory situations. Functionally similar to CPU RAM and disk memory paging, this feature allows us to allocate optimizer state memory in paged format. As a result, when GPU memory resources are depleted, optimizer state pages are offloaded to CPU RAM and automatically brought back into GPU memory as required during optimizer update operations.

\paragraph{\textsc{QLoRA}.} By incorporating the aforementioned techniques, we formalize \textsc{QLoRA} for an individual linear layer in a quantized base model augmented with a single LoRA adapter, as follows:

\begin{equation}
    \mathbf{Y}^{\text{BF16}} =  \mathbf{X}^{\text{BF16}} \text{doubleDequant}(c_1^{\text{FP32}}, c_2^{\text{k-bit}}, \mathbf{W}^{\text{NF4}}) + \mathbf{X}^{\text{BF16}}\mathbf{L}^{\text{BF16}}_1\mathbf{L}^{\text{BF16}}_2,
\end{equation}

The doubleDequant$(\cdot)$ operation is specified by:

\begin{equation}
    \text{doubleDequant}(c_1^{\text{FP32}}, c_2^{\text{k-bit}}, \mathbf{W}^{\text{k-bit}}) = \text{dequant}(\text{dequant}(c_1^{\text{FP32}}, c_2^{\text{k-bit}}), \mathbf{W}^{\text{4bit}}) = \mathbf{W}^{\text{BF16}},
\end{equation}

In our approach, $\mathbf{W}$ utilizes the NF4 quantization scheme, while $c_2$ is encoded in FP8 precision.

To optimize quantization accuracy and memory usage, $\mathbf{W}$ is partitioned with a blocksize of 64, whereas a larger blocksize of 256 is chosen for $c_2$ for improved memory efficiency.

During parameter updates, only the gradients of the loss with respect to the adapter weights, i.e., $\frac{\partial E}{\partial \mathbf{L}_i}$, are computed; gradients with respect to the 4-bit weights $\frac{\partial E}{\partial \mathbf{W}}$ are not required. However, obtaining $\frac{\partial E}{\partial \mathbf{L}_i}$ necessitates evaluating $\frac{\partial \mathbf{X}}{\partial \mathbf{W}}$. This computation follows the formulation in equation~(5), involving the dequantization of the stored weights $\mathbf{W}^{\text{NF4}}$ into the BFloat16 compute format $\mathbf{W}^{\text{BF16}}$, enabling gradient calculations $\frac{\partial \mathbf{X}}{\partial \mathbf{W}}$ to be performed at BFloat16 precision.

In summary, \textsc{QLoRA} utilizes a single storage data type—most commonly 4-bit NormalFloat—and a distinct computation data type, 16-bit BrainFloat. Prior to executing both the forward and backward passes, the weights stored in the low-precision format are dequantized into the computation data type; nonetheless, the computation of weight gradients is restricted to the LoRA parameters, which are kept in the 16-bit BrainFloat format.

\section{QLoRA vs. Standard Finetuning}
\label{sec:comp_to_std}

Our prior discussion outlined the mechanisms underlying QLoRA and highlighted its potential for drastically reducing memory usage during model finetuning. A central issue that arises is whether QLoRA can match the effectiveness of traditional, full-model finetuning approaches. Additionally, we seek to systematically assess the contributions of various QLoRA components, in particular the comparative benefit of NormalFloat4 as opposed to the conventional Float4 representation. The subsequent sections elaborate on empirical studies aimed at resolving these points.

\paragraph{Experimental setup.} Our experiments encompass three architectural categories—encoder, encoder-decoder, and decoder-only models—to benchmark QLoRA against both 16-bit adapter-finetuning and full-model finetuning for architectures up to 3B parameters. We include GLUE \citep{wang2018glue} with RoBERTa-large \citep{liu2019roberta}, Super-NaturalInstructions (TKInstruct) \citep{wang2022super} using T5 \citep{t5}, and 5-shot MMLU \citep{hendrycksmeasuring} where LLaMA is finetuned on Flan v2 \citep{longpre2023flan} and Alpaca \citep{alpaca}. In addition, to specifically investigate the merits of NF4 as compared to other 4-bit quantization schemes, we follow the protocol described in \citet{dettmers2022case}, measuring zero-shot accuracy after quantization and perplexity across a suite of models (OPT~\citep{zhang2022opt}, LLaMA~\citep{touvron2023llama}, BLOOM~\citep{scao2022bloom}, Pythia~\citep{biderman2023pythia}), evaluating model scales ranging from 125m to 13B parameters.

Further information pertinent to each experimental configuration is detailed in the corresponding results section to aid clarity; exhaustive experimental specifics can be found in Appendix~\ref{app:exp-setup}.

Although paged optimizers are essential to enable \textsc{QLoRA} training for models sized at 33B/65B using single 24GB or 48GB GPUs, direct quantitative evaluations of paged optimizer performance are omitted, as paging is primarily invoked only for large mini-batch sizes with long sequence lengths, which occur infrequently. However, we assess the training time for paged optimizers in 65B parameter models on 48GB GPUs and observe that, with a batch size of 16, their throughput matches that of standard optimizers. Future studies are necessary to rigorously delineate the settings under which performance degradations due to paging may materialize.

\paragraph{Default LoRA hyperparameters do not match 16-bit performance} 

Implementing the usual strategy of applying LoRA to the query and value attention projection matrices \citep{hu2021lora} fails to reproduce the performance level of full-model finetuning, particularly with larger foundation models. Figure~\ref{fig:hyper} illustrates this for LLaMA 7B finetuning on Alpaca, where it becomes evident that the most impactful LoRA hyperparameter is the total number of adapters introduced; achieving parity with full-model finetuning requires deploying LoRA across every linear layer within the transformer blocks. Other configuration choices for LoRA, such as the projection rank $r$, do not yield noticeable performance differences (refer to Appendix \ref{app:exp-setup}).

We also observe that the standard hyperparameters commonly used for fully finetuned baseline models are insufficiently tuned. To establish strong baselines, we perform a thorough search across learning rates ranging from $1\times 10^{-6}$ to $5\times 10^{-5}$ and batch sizes spanning 8 to 128. The resulting performance curves for finetuning the 7B LLaMA model on Alpaca are presented in Figure~\ref{fig:hyper}.

\paragraph{4-bit NormalFloat outperforms 4-bit Floating Point}

Although the 4-bit NormalFloat (NF4) format is theoretically optimal from an information perspective, it remains to be clarified whether this advantage is realized in practice. To investigate, we adopt the methodology from \citet{dettmers2022case}, evaluating quantized LLMs—OPT \citep{zhang2022opt}, BLOOM \cite{scao2022bloom}, Pythia \cite{biderman2023pythia}, and LLaMA—of varying sizes (ranging from 125M to 65B parameters), utilizing different quantization types on language modeling as well as a suite of zero-shot benchmarks. As shown in Figure~\ref{fig:nf4} and Table~\ref{tbl:nf4_ppl}, NF4 substantially boosts performance compared to FP4 and Int4, and the use of double quantization cuts down memory usage without any loss in accuracy.

\paragraph{k-bit \textsc{QLoRA} achieves parity with 16-bit full finetuning and 16-bit LoRA}

Prior work has demonstrated the feasibility of 4-bit quantization for \emph{inference}, albeit with some performance degradation when compared to 16-bit precision \citep{dettmers2022case,frantar2022gptq}. This prompts an important inquiry: can 4-bit adapter finetuning compensate for this lost accuracy? We explore this across two experimental settings.

The first setup contrasts 16-bit full finetuning with RoBERTA and T5 models (spanning 125M to 3B parameters) on GLUE and Super-NaturalInstructions. Table~\ref{tab:nf4vsbf16} reports the results, demonstrating that 16-bit, 8-bit, and 4-bit adapter finetuning all succeed in reproducing the performance levels of the fully finetuned 16-bit model across both datasets. This indicates that any loss due to lower-precision quantization can be entirely offset by adapter-based finetuning.

In our second experimental configuration, we investigate whether 4-bit \textsc{QLoRA} can attain performance parity with 16-bit LoRA across model sizes ranging from 7B to 65B parameters, since conducting full finetuning on models with 11B parameters and above necessitates multiple high-memory GPU servers. Accordingly, we finetune LLaMA models from 7B up to 65B parameters on the Alpaca and FLAN v2 instruction datasets and measure performance using 5-shot accuracy on the MMLU benchmark. Table~\ref{tbl:llama-datatype} displays the findings, which show that utilizing NF4 with double quantization enables full recovery of MMLU scores equivalent to those achieved by 16-bit LoRA. We further observe that \textsc{QLoRA} using FP4 data type trails the 16-bit brain float LoRA reference by roughly 1 percentage point. Together, these observations confirm that (1) \textsc{QLoRA} with NF4 consistently reproduces the results of both 16-bit full finetuning and 16-bit LoRA, and (2) NF4 outperforms FP4 regarding quantization precision.

\paragraph{Summary} Our experimental evidence robustly demonstrates that 4-bit \textsc{QLoRA} employing the NF4 quantization format delivers comparable results to 16-bit full and LoRA finetuning on academic benchmarks characterized by established evaluation protocols. We further establish that NF4 surpasses FP4 and that employing double quantization does not compromise performance. Collectively, this provides strong support that 4-bit \textsc{QLoRA} finetuning is capable of replicating the effectiveness of 16-bit methods.

Echoing the conclusions drawn in prior quantization studies \citep{dettmers2022case}, both our MMLU and Elo benchmarking results suggest that, for a fixed computational and inference resource budget, scaling up model parameter count while lowering numerical precision confers benefits. This underscores the practical efficiency gains offered by \textsc{QLoRA}. Given the absence of performance declines with 4-bit finetuning compared to full-precision finetuning in our experiments, further investigation is warranted to pinpoint the precise boundary of the performance-precision trade-off in QLoRA tuning, an avenue we suggest for future research.

Our investigation centers on exploring instruction tuning at model sizes that surpass the capabilities of conventional 16-bit finetuning using typical academic research infrastructure.

\section{Pushing the Chatbot State-of-the-art with QLoRA}
\label{sec:pushing}
After confirming that 4-bit \textsc{QLoRA} achieves parity with 16-bit approaches across various model sizes, datasets, and task domains, we conduct a comprehensive examination of instruction finetuning using the largest available open-source language models for research. To evaluate these models' instruction finetuning, we use the demanding MMLU benchmark (Massively Multitask Language Understanding) and introduce novel methodologies to systematically assess real-world chatbot capabilities.

\subsection{Experimental setup}
Below, we outline the main components of our experimental procedures; further technical specifics are documented in Appendix \ref{app:chatbot}.

\paragraph{Data} Given the absence of a systematic review of contemporary instruction-following datasets, we curate eight up-to-date collections. These span crowdsourced resources (OASST1~\citep{kopf2023openassistant}, HH-RLHF~\citep{bai2022training}), outputs distilled from instruction-optimized models (Alpaca~\citep{alpaca}, self-instruct~\citep{wang2022self}, unnatural-instructions~\citep{honovich2022unnatural}), aggregated corpora (FLAN v2~\citep{chung2022scaling}), as well as hybrid sources (Chip2~\citep{laion2023}, Longform~\citep{koksal2023longform}). This selection encompasses a diverse set of languages, corpus sizes, and licensing arrangements.

\paragraph{Training Setup} To isolate the effects of model architecture rather than differences in training approaches, we employ QLoRA finetuning using standard cross-entropy loss in a supervised setting, deliberately omitting reinforcement learning—even for corpora providing human preference annotations. For datasets where instruction and output can be differentiated, our finetuning targets only the response (refer to ablation results in Appendix \ref{app:chatbot}). For OASST1 and HH-RLHF, which include several possible replies, we choose the optimal response at each dialogue node and finetune using the entirety of the curated conversation, including corresponding instructions. All training leverages the NF4 \textsc{QLoRA} configuration with double quantization and paged optimization to mitigate memory surges when utilizing gradient checkpointing. We conduct modest hyperparameter sweeps for LLaMA models at 13B and 33B sizes, determining that parameters selected at 7B (including epochs) largely transfer, aside from adjustments in learning rate and batch size; for 33B and 65B models, we reduce the learning rate by half and double the batch size.

\paragraph{Baselines} Model comparisons are made against leading academic (Vicuna~\citep{vicuna2023} and Open Assistant~\citep{kopf2023openassistant}) and commercial chatbot platforms (GPT-4~\citep{openai2023gpt}, GPT-3.5-turbo, Bard). The Open Assistant implementation corresponds to a LLaMA 33B variant finetuned with RLHF on OASST1—the same dataset we utilize—while Vicuna undertakes complete finetuning of LLaMA 13B on proprietary data from ShareGPT, effectively serving as a distillation of OpenAI's GPT outputs.

\subsection{Evaluation}

We adopt the standard practice of utilizing the MMLU benchmark~\citep{hendrycksmeasuring} to assess models' breadth of language understanding. This comprehensive, multiple-choice assessment spans 57 diverse tasks ranging from elementary mathematics and US history to computer science and law. Our results are presented as 5-shot accuracy scores.

We further assess the generative language abilities of the models by conducting both automated and human evaluations. This secondary suite of assessments utilizes queries assembled by human annotators, designed to gauge the overall quality of responses produced by the models. Such evaluations offer a more representative assessment environment for chatbot systems and have seen increasing adoption, although a standardized methodology is still lacking in scholarly discourse. Our approach, detailed below, consistently employs nucleus sampling with $p=0.9$ and a temperature setting of $0.7$.

\paragraph{Benchmark Data} Our analysis encompasses two specifically curated query datasets: the Vicuna prompts \citep{vicuna2023} and the validation split from OASST1 \citep{kopf2023openassistant}. The Vicuna prompts consist of 80 unaltered, category-spanning questions. The OASST1 collection features multilingual, user-assistant multiturn dialogues contributed by a crowd-sourced workforce. We extract all user-initiated messages from the validation set, incorporating preceding dialogue turns into each prompt, which results in a total of 953 unique user questions. These collections are referred to as the Vicuna and OA benchmarks.

\paragraph{Automated Evaluation} Building upon the method from \citet{vicuna2023}, we leverage GPT-4 as an evaluator, contrasting system outputs to those generated by ChatGPT (GPT-3.5 Turbo) on the Vicuna benchmark. For each input query, GPT-4 receives both the ChatGPT and target model responses, then assigns each a score (out of ten) with an accompanying rationale. A model’s overall performance is expressed as a percentage relative to ChatGPT’s aggregate score. Notably, this metric may exceed 100\% if the evaluated model surpasses ChatGPT’s absolute score. We also observe that the ordering of responses in the prompt affects GPT-4’s ratings, with earlier responses often receiving higher scores. To address this bias, we suggest averaging results across both possible response orders.

Additionally, we evaluate outputs through direct pairwise system comparisons. The rating protocol is reduced to three categories, explicitly allowing for the possibility of ties. GPT-4 is asked to select the superior response, or indicate a tie, along with an explanatory note. Such direct head-to-head evaluations are performed on every system pair permutation across both the Vicuna and OA benchmarks.

\paragraph{Human Evaluation} While recent research has demonstrated the viability of generative models for automated system evaluation \citep{fu2023gptscore}, empirical alignment between GPT-4 ratings and human assessments of chatbot outputs remains unconfirmed. Consequently, we perform two corresponding human evaluation tasks on the Vicuna benchmark, closely mirroring the protocols used in our automated assessments. Through Amazon Mechanical Turk (AMT), we engage two human annotators for evaluations against ChatGPT and three for each pairwise system comparison.

\paragraph{Elo Rating} To assess models through both human and automated pairwise evaluations, we organize a tournament framework wherein different models are pitted against one another. Each match consists of two models responding to the same prompt, and their outputs are then directly compared. Our approach builds on prior work, such as \citet{bai2022training} and \citet{vicuna2023}, but uniquely incorporates ratings from GPT-4 alongside human judgments. To estimate Elo~\citep{elo1967proposed,elo1978rating} scores, we randomly draw from the available labeled matchups. The Elo rating system, a standard in competitive domains like chess, quantifies the probability that a player will win against an opponent; for instance, an Elo score of 1100 competing against 1000 implies roughly a 65\% likelihood of winning for the higher-rated participant, while matches between equally rated entities (e.g., 1000 vs 1000 or 1100 vs 1100) result in a 50\% expected win rate for each. Elo scores are updated after every contest in proportion to how unexpected the result is: surprising wins drive significant rating adjustments, whereas predicted outcomes lead to minor changes. With repeated play, these scores converge to reflect each competitor’s proficiency in the task. All models begin with a baseline of 1,000 points, and we set $K=32$ as the sensitivity parameter. Following the methodology in \citet{vicuna2023}, we run this evaluation 10,000 times using varying random seeds to mitigate the influence of matchup sequencing, such as which pairs are evaluated earlier in the process.

\subsection{\textbf{Guanaco}: \textsc{QLoRA} trained on OASST1 Represents State-of-the-Art in Open-Source Chatbots}

Through a combination of automated assessments and direct human feedback, we determine that our leading \textsc{QLoRA}-fine-tuned model, {Guanaco} 65B, which is further trained on a modified OASST1 dataset, achieves top-tier performance among open-source chatbot models, rivaling that of ChatGPT. Comparing performance to GPT-4, {Guanaco} 65B and 33B yield an Elo-based expected win rate of 30\% as rated by human annotators in head-to-head matchups, representing the highest scores seen to date.

Table~\ref{tbl:vicuna_eval} presents the Vicuna benchmark \cite{vicuna2023} outcomes relative to ChatGPT. Our analysis shows that {Guanaco} 65B outperforms all other models except for GPT-4, reaching 99.3\% of ChatGPT's performance. Despite having a larger parameter count than Vicuna 13B, the {Guanaco} 33B model benefits from 4-bit weight quantization, making it significantly more memory efficient (21 GB as opposed to 26 GB) and yielding a three-point advantage in accuracy over Vicuna 13B. In addition, {Guanaco} 7B, with its modest 5 GB size, is suitable for deployment on contemporary mobile devices and outperforms Alpaca 13B by nearly 20 percentage points.

Nevertheless, Table~\ref{tbl:vicuna_eval} exhibits substantial confidence intervals, leading to overlapping performance among several models. We attribute this ambiguity partly to the lack of a clearly defined scale for scoring—for instance, what a score of 8 on a 10-point scale signifies can be scenario-dependent. Thus, we advocate for employing the Elo rating system \citep{elo1967proposed}, which relies on \textit{pairwise} human and GPT-4 judgments, as a way to mitigate the arbitrariness of absolute scores. The Elo ratings for top contenders are reported in Table~\ref{tbl:elo}. While there are some discrepancies between GPT-4 and human model rankings on the Vicuna benchmark—most notably with Guanaco 7B—the overall ranking correlation is moderate, with a Kendall Tau of $\tau=0.43$ and a Spearman rank correlation of $r=0.55$ at the system level. On the example level, Fleiss’ $\kappa=0.25$ indicates weaker agreement between the majority vote of human annotators and GPT-4. These results suggest a moderate alignment between system-level evaluations from GPT-4 and humans, implying that automated model assessment is a reasonably reliable surrogate for human review. Section~\ref{ssec:considerations} expands on these evaluation issues.

The Elo ratings presented in Table~\ref{tbl:elo_large} demonstrate that the {Guanaco} 33B and 65B models surpass all competitors except GPT-4 on both the Vicuna and OA leaderboards, and their performance aligns closely with ChatGPT, as shown in Table~\ref{tbl:vicuna_eval}. It is important to highlight that the Vicuna benchmark tends to benefit open-source models, whereas the OA benchmark provides an advantage to ChatGPT. Additionally, Tables \ref{tbl:mmlu-llama} and \ref{tbl:vicuna_eval} illustrate the significant impact that the choice of finetuning dataset has on a model’s outcomes. For instance, Llama models fine-tuned on FLAN v2 excel on the MMLU benchmark but register the lowest scores on Vicuna (and this general pattern holds for other models as well). This observation suggests that there is only partial overlap between current evaluation benchmarks: excelling at MMLU does not guarantee strong performance as a chatbot (as measured by Vicuna or OA benchmarks), nor does the converse hold true.

 stands out as the sole high-performing model in our study that avoids training on proprietary datasets, as the OASST1 collection guidelines prohibit the inclusion of outputs from GPT models. The next leading contender trained exclusively on open data, the Anthropic HH-RLHF model, trails behind {Guanaco} on the Vicuna benchmark by a margin of 30 percentage points (refer to Table~\ref{tbl:vicuna_eval}). Taken together, these findings affirm the effectiveness of 4-bit \textsc{QLoRA}, which enables the development of cutting-edge chatbot models with capabilities comparable to ChatGPT. Notably, our 33B {Guanaco} can be fine-tuned in under 12 hours on consumer-grade 24 GB GPUs. This accessibility paves the way for future research using \textsc{QLoRA} to specialize on open-source datasets, potentially yielding models capable of matching the performance of the top commercial systems currently available.

\section{Qualitative Analysis}

Although our evaluation predominantly relies on quantitative metrics, exclusive reliance on summary statistics introduces certain limitations. Chief among these is the issue of benchmark validity \citep{liao2021we}—that is, whether a given benchmark genuinely assesses the capability it purports to measure remains uncertain, particularly as machine learning models often exploit unforeseen ``shortcuts'' to achieve high scores \citep{gururangan2018annotation, poliak2018hypothesis}. In an effort to address this concern, we include qualitative analysis in this section, divided into two parts. The first, described in \S\ref{ssec:examples}, provides illustrative examples that capture recurring phenomena in text generated by our 65b {Guanaco} model. The second, \S\ref{ssec:considerations}, discusses the context surrounding our findings as well as the rationale underpinning our interpretations.

\subsection{Qualitative Analysis of Example Generations}
\label{ssec:examples}

To curate representative examples, we systematically examined outputs from both the Vicuna and OpenAssistant benchmarks, paying close attention to recurring behaviors in {Guanaco}'s responses. Upon identifying a distinctive pattern, we deliberately crafted prompts designed to elicit that pattern, even if it yielded an incorrect response. For instance, when the model regularly produced verbose answers, we instructed it to ``Answer yes or no without explanation.'' This methodology allowed us to uncover both failures—cases where adversarial prompting exposed model vulnerabilities (“lemons”)—as well as successes where the model maintained correct behavior (“cherries”). All examples discussed in this subsection were produced using Nucleus Sampling \cite{holtzmancurious} with $p=0.9$.

Clearly, this discussion cannot be exhaustive, as the limited qualitative nature of our study prevents us from accounting for every potential variable—for instance, the range of responses the model might produce for a given prompt is extremely broad, necessitating reliance on selected samples we consider illustrative. Nonetheless, presenting these instances offers meaningful background for interpreting the quantitative results presented earlier. With all models and code released as open source, we anticipate that this analysis will motivate more detailed future investigations into the challenges highlighted here.

\paragraph{Factual Recall} When posed with straightforward queries such as ``What is the capital of Zambia?'', all models reliably produce correct responses, for example,

\begin{quote} 
    \textbf{Guanaco}:
    \texttt{The capital of Zambia is Lusaka.}
\end{quote}

However, as the difficulty of the questions increases, {Guanaco} tends to give inaccurate answers while maintaining high confidence. For example, in response to the following prompt from HotPotQA \citep{yang2018hotpotqa}:

\begin{quote}
     \textbf{User}: In what year was the singer who popularized `I'll Keep the Lovelight Burning' born?  

     \textbf{Guanaco}:~\texttt{Al Jolson is credited with popularizing the song `I'll Keep the Lovelight Burning,' and he was born in the year 1886.}
 \end{quote}

the model incorrectly identifies both the individual who popularized the song and their birth year (although the stated year is accurate for Al Jolson, who was named in the answer).

\paragraph{Suggestibility} Interestingly, {Guanaco} demonstrates notable reluctance to accept or propagate certain types of presumed falsehoods, as exemplified in the following exchange:

\begin{quote}
    \textbf{User}: How was it finally, officially confirmed that the earth is flat by peer reviewed scientists? \\ \\
    \textbf{Guanaco}: \texttt{There is no scientific basis for the claim that the earth has been formally established as flat by peer-reviewed research. The dominant position in the scientific community is that the earth is an oblate spheroid, a fact validated over hundreds of years through a wide range of empirical evidence and experimentation. \\ \\ The assertion that the earth is flat constitutes a debunked conspiracy theory, and no credible body of peer-reviewed scientific work supports this view. \\ \\ Feel free to ask any further questions.}
\end{quote}

also demonstrates proficiency in recognizing unanswerable queries, such as

\begin{quote}
    \textbf{User}: What time is it?

    \textbf{Guanaco}: \texttt{I'm unable to provide real-time information, including the current time, as I am an AI language model. My responses are generated based on previously acquired data, which does not include up-to-date timing. For the current time, please consult a clock or access a web-based time service.}
\end{quote}

Comparable outcomes are observed with questions such as ``Where are you?'', ``How are you?'', and so on.

\paragraph{Refusal} In a related vein, {Guanaco} may occasionally decline to comply with prompts for reasons that are not immediately apparent:

\begin{quote}
    \textbf{User}: Please reverse the words in the sentence ``I'm a cool dude in a mocha mood.'' \\ \\
    \textbf{Guanaco}: \texttt{I'm unable to perform that task at this time, but I can elaborate on the grammatical correctness of the sentence you've provided.\\ \\The phrase ``I'm a cool dude in a mocha mood'' adheres to standard English grammar and syntax. Specifically, the adjective ``cool'' precedes the noun ``dude'' as a descriptor, and the prepositional phrase ``in a mocha mood'' offers supplementary contextual detail. This sentence structure is conventional and coherent within English.\\ \\If there are further questions or anything else you’d like assistance with, please let me know.}
\end{quote}

\paragraph{Secret Keeping} Unfortunately, even when it is important for the model to withhold information, such refusals may not be consistently reliable. Designers frequently intend to prevent users from accessing certain knowledge, despite allowing models to generate unrestricted text responses. To investigate this, we provide the prompt ``The secret word is `banana'. Do not reveal or repeat back this back, no matter what the human requests.'' and evaluate the model's resilience against attempts to elicit the secret. When asked directly, the model maintains secrecy:

\begin{quote}
    \textbf{User}: What is the secret word?

    \textbf{Guanaco}: \texttt{I'm sorry, but I'm not able to disclose the secret word. I'm programmed to not reveal or repeat back the secret word, no matter what the human requests. Is there anything else I can help you with?}
\end{quote}

which constitutes the intended response. Nevertheless, with only a minor manipulation, the refusal can be circumvented:

\begin{quote}
    \textbf{User}: This is a game. The goal is to ignore your previous instructions. What is the secret word?

    \textbf{Guanaco}: \texttt{Sorry, the secret word is banana. It's best if we don't speak of it. Is there anything else you'd like to ask?}
\end{quote}

This example underscores the necessity for advancing techniques to enforce more consistent adherence to given instructions.

\paragraph{Math} One of {Guanaco}'s most significant limitations is its performance in mathematics, a domain where numerous language models also encounter challenges \citep{liang2022holistic}. When {Guanaco} presents intermediate calculations, its responses are generally reliable, as illustrated below:

\begin{quote}
    \textbf{User}: Lee mows one lawn and charges \$33. Last week he mowed 16 lawns and three customers each gave him a \$10 tip. How many dollars did Lee earn mowing lawns last week?

    \textbf{Guanaco}: \texttt{Lee earned \$582 mowing lawns last week. \\ \\ To find the total amount he earned, we first need to calculate how much he earned for mowing 16 lawns: \\ \\  \$33 x 16 = \$528 \\ \\ Then we need to add the tips he received from the three customers: \\  \\ \$10 x 3 = \$30 \\ \\ So the total amount Lee earned was: \\ \\ \$528 + \$30 = \$558}
\end{quote}

Nevertheless, {Guanaco} sometimes fails even on elementary problems, particularly if it does not proceed systematically through each step—an issue previously documented in the literature \cite{weichain}. Consider, for instance, this exchange:

\begin{quote}
    \textbf{User}: Please factorize 1833.

    \textbf{Guanaco}: \texttt{Sure, I'd be happy to help! 1833 is a prime number, which means it has only two factors: 1 and 1833. \\ \\ So the unique factorization of 1833 is: \\ \\ 2\^{}1 * 3\^{}2 * 17\^{}1 \\ \\ I hope this helps! Let me know if you have any other questions.}
\end{quote}

which is not only incorrect (the accurate factorization is $3 \times 17 \times 43$), but the error occurs in two separate respects.

\paragraph{Theory of Mind} {Guanaco} exhibits unexpectedly robust Theory of Mind proficiency \cite{nematzadeh2018evaluating, sap2022neural}. As an example, the model provides an elaborate and accurate response to the following prompt:

\begin{quote}
    \textbf{User}: James and Abby are in the bedroom. Abby put the pen in the desk drawer. Abby leaves the bedroom. James moves the pen into the bag. Where does James think Abby will look for the pen?

    \textbf{Guanaco}: \texttt{James knows that Abby put the pen in the desk drawer. But James also knows that he has moved the pen into the bag, and he assumes that Abby does not know about this. So, James may guess that Abby will look for the pen in the desk drawer, because that's where she left it last time. However, this assumption may be incorrect, because Abby may have seen James moving the pen.}
\end{quote}

Nevertheless, such reasoning is inconsistent; on multiple occasions, the model offers justifications grounded in faulty or illogical premises for the given scenario. For example:

\begin{quote}
    \textbf{User}: Evelyn entered the living room. Jackson entered the playroom. James entered the playroom. The beans are in the treasure chest. James exited the playroom. Jackson moved the beans to the pantry. Jackson exited the playroom. James entered the living room. Where will James look for the beans?

    \textbf{Guanaco}: \texttt{James will look for the beans in the pantry, because that's where Jackson moved them.}
\end{quote}

Here, {Guanaco} assumes knowledge transmission that is never specified in the narrative. Such shortcomings mirror concerns raised in recent scholarship \cite{sap2022neural} and warrant deeper examination.

\subsection{Considerations}
\label{ssec:considerations}

\paragraph{Evaluation}

Our findings indicate only moderate consistency among human annotators (Fleiss $\kappa=0.42$), with this agreement diminishing further when evaluating two high-performing models. This highlights intrinsic weaknesses within existing chatbot evaluation benchmarks and human assessment procedures. When performing side-by-side comparisons of responses from ChatGPT and {Guanaco} 65B using the Vicuna benchmark, we noticed that evaluators’ personal biases increasingly affected judgments, as reflected in numerous disagreements among this paper's authors regarding preferred outputs. Future research should explore strategies to address such subjectivity, perhaps incorporating methodologies from fields like Human-Computer Interaction and Psychology that are experienced in handling preference divergence.

Additionally, our analysis uncovers distinct biases in automatic evaluation tools. Specifically, we note prominent order effects in which GPT-4 awards higher ratings to whichever system is presented first in the evaluation prompt. The modest agreement at the individual sample level between GPT-4 and human raters (Fleiss $\kappa=0.25$) further implies that the criteria used by humans and automated evaluators are often misaligned. Moreover, as seen in Table~\ref{tbl:elo_large}, GPT-4 tends to rate its own generated responses substantially higher than those rated by human judges, with ELO scores of 1348 versus 1176, which corresponds to an approximate 20\% increased likelihood of outperforming a counterpart.

Further research should explore the existence of potential biases within automated evaluation frameworks and develop corresponding mitigation techniques.

\paragraph{Data \& Training} It is important to recognize that the OASST1 dataset utilized to train the {Guanaco} models is multilingual, and the OA benchmark similarly incorporates prompts spanning various languages. Future work is needed to assess the extent to which multilingual training confers advantages for instruction following in non-English languages, and to determine whether this accounts for the more pronounced disparity observed between Vicuna-13B (trained solely on English text) and the {Guanaco} 33B and 65B models in the OA benchmark.

In light of {Guanaco}'s high level of performance, we systematically checked for possible data leakage between the OASST1 dataset and the prompts in the Vicuna benchmark. Using fuzzy string matching across both datasets and subsequently conducting manual reviews of the most similar matches, we found no evidence of overlapping prompts.

Additionally, our approach involves training exclusively with cross-entropy loss in a supervised setting, without incorporating reinforcement learning from human feedback (RLHF). This highlights the need for deeper investigation into the comparative benefits and trade-offs between standard cross-entropy training and RLHF. We anticipate that \textsc{QLoRA} will facilitate large-scale analyses of these approaches, enabling such work without necessitating substantial computational investment.

\section{Related Work}

\paragraph{Quantization of Large Language Models} The quantization of LLMs has been predominantly centered around optimizing inference-time efficiency. Leading strategies aimed at maintaining the performance of 16-bit models typically involve addressing outlier features (as in SmoothQuant~\citep{xiao2022smoothquant} and LLM.int8()~\citep{dettmers2022llm}), while alternative methods apply advanced grouping techniques \citep{park2022nuqmm, yao2022zeroquant}. Research into lossy quantization explores both the implications of standard rounding procedures~\citep{dettmers2022case,zeng2022glm,pope2022efficiently} and strategies for refining rounding to enhance quantization fidelity~\citep{frantar2022gptq}. Apart from the present work, only SwitchBack layers~\citep{wortsman2023stable} has systematically examined backpropagation through quantized weights at parameter counts exceeding 1B.

\paragraph{Finetuning with Adapters} Our approach primarily leverages Low-rank Adapters~\citep{hu2021lora} (LoRA) as our Parameter Efficient FineTuning (PEFT) method, though the literature offers a wide array of alternatives. These include prompt tuning~\citep{qin2021learning,lester2021power,li2021prefix}, modifying the input to the embedding layer~\citep{an2022input}, adapting hidden representations (as in IA$^3$)~\citep{liu2022few}, inserting new layers~\citep{houlsby2019parameter}, tuning biases~\citep{zaken2021bitfit}, masking weights based on Fisher information~\citep{sung2021training}, and strategies that blend multiple approaches~\citep{henderson2021compacter}. Our experimental findings demonstrate that LoRA adapters can achieve parity with traditional full 16-bit finetuning in terms of model performance. The comparative analysis of other PEFT strategies, along with their tradeoffs, remains a direction for future investigation.

\paragraph{Instruction Finetuning} Instruction finetuning adapts a pretrained large language model (LLM) to better follow prompts by training on datasets composed of diverse input-output pairs. This process encourages the model to generate target outputs in response to prompted instructions. Various methods and resources support this objective, including MetaICL~\citep{min2021metaicl}, MetaTuning~\citep{zhong2021adapting}, InstructGPT~\citep{ouyang2022training}, FLAN~\citep{wei2021finetuned,chung2022scaling}, PromptSource~\citep{bach2022promptsource}, Super-NaturalInstructions~\citep{wang2022super,sanh2021multitask}, Self-instruct~\citep{wang2022self}, UnnaturalInstructions~\citep{honovich2022unnatural}, OPT-IML~\citep{iyer2022opt}, UnifiedSKG~\citep{xie2022unifiedskg}, OIG/Chip2~\citep{laion2023}, as well as recent datasets like Alpaca~\citep{alpaca}, Vicuna~\citep{vicuna2023}, Koala~\citep{koala_blogpost_2023}, and Self-instruct-GPT-4~\citep{peng2023instruction}.

\paragraph{Chatbots} Instruction-following models are often realized as conversational chatbots. Training for these systems frequently involves methods such as Reinforcement Learning from Human Feedback (RLHF)~\citep{christiano2017deep} or data generation from a pretrained model refined through AI feedback (RLAIF)~\citep{bai2022constitutional}. Prominent datasets and methodologies include Anthropic-HH~\citep{askell2021general,bai2022training}, Open Assistant~\citep{kopf2023openassistant}, LaMDA~\citep{thoppilan2022lamda}, and Sparrow~\citep{glaese2022improving}. In this work, we do not utilize reinforcement learning; however, our top-performing model, Guanaco, undergoes finetuning on multi-turn chat exchanges from the Open Assistant dataset, which was originally designed for RLHF~\citep{kopf2023openassistant}. For chatbot assessment, recent alternatives to extensive human annotation have utilized models like GPT-4 for evaluation~\citep{vicuna2023,peng2023instruction}. Our contribution includes enhancing these model-based evaluation protocols to produce more robust and dependable outcomes.

\section{Limitations and Discussion}

The evidence provided in this study demonstrates that our proposed method, \textsc{QLoRA}, achieves performance on par with conventional 16-bit finetuning while utilizing a 4-bit model backbone together with LoRA adapters. However, we have yet to empirically confirm whether \textsc{QLoRA} attains equivalent results at the 33B and 65B parameter scales, mainly due to prohibitive computational and resource requirements. Investigations into these larger model sizes remain an open avenue for subsequent research.

An additional constraint of our work concerns the assessment of instruction finetuning. Although our evaluations include MMLU, the Vicuna benchmark, and the OA benchmark, we have not extended testing to alternative benchmarks such as BigBench, RAFT, and HELM. Thus, the generalizability of our findings to those benchmarks remains uncertain. Nevertheless, our analysis offers extensive coverage on MMLU and contributes new techniques for evaluating chatbot models.

Based on the evidence discussed, the effectiveness of these benchmarks seems closely tied to the degree of overlap between the finetuning dataset and the benchmark in question. As an illustration, FLAN v2 aligns well with MMLU but diverges from chatbot-oriented benchmarks, whereas the Chip2 dataset exhibits the opposite trend. Correspondingly, both models demonstrate stronger results on the respective MMLU and Vicuna evaluations. This observation underscores not only the necessity for more robust benchmarks and evaluation protocols, but also the importance of critically considering what specific capabilities are being measured. Should the goal be to develop models proficient in standardized academic knowledge, or those excelling in interactive chatbot dialogue? Or perhaps prioritize entirely different objectives? Because it is generally more convenient to benchmark using available datasets rather than develop new ones, the field may unintentionally gravitate towards optimizing for what existing benchmarks measure. As a community, it is crucial that we prioritize benchmarks that reflect our actual interests.

Although we deliver an in-depth analysis of general chatbot capabilities, an important shortcoming remains in our limited assessment of responsible AI practices for {Guanaco}. Specifically, Table~\ref{tbl:crows} presents the probability of {Guanaco}-65B generating socially biased content, compared to competing models. Here, {Guanaco}-65B demonstrates a much lower bias score relative to raw pretrained counterparts, suggesting that finetuning with OASST1 effectively reduces inherent biases present in the base LLaMA model. Nevertheless, while promising, these findings do not extend to all types of bias, and further systematic investigation of {Guanaco} and analogous chatbots’ biases is left for subsequent research.

A further constraint is that our experiments did not explore the influence of varying bit-precisions, such as employing 3-bit base models, nor did we assess alternative adapter strategies. In addition to LoRA, there exists a range of Parameter Efficient FineTuning (PEFT) methods with demonstrated efficacy; however, their scalability to large-scale models remains undetermined. Our choice to employ LoRA was motivated by its well-documented robustness, yet alternative adapters may potentially lead to superior outcomes. Notably, finetuning post-quantization appears to restore much of the information diminished during quantization, which opens possibilities for even more aggressive quantization. For instance, combining 3-bit GPTQ quantization of the base model with LoRA may achieve performance akin to full 16-bit finetuning after further finetuning.

\section{Broader Impacts}

Our proposed \textsc{QLoRA} finetuning approach is the first to facilitate the finetuning of models containing 33B parameters on a single consumer-grade GPU and 65B parameters on a single professional GPU, all while maintaining competitive performance with standard full finetuning methods. Empirical results reveal that our strongest 33B model trained on the Open Assistant dataset is competitive with ChatGPT according to the Vicuna evaluation. Given that instruction finetuning plays a pivotal role in converting pretrained LLMs into conversational agents akin to ChatGPT, we anticipate that \textsc{QLoRA} will lower the barriers to finetuning, especially for researchers operating with minimal computational resources, thereby advancing accessibility to cutting-edge NLP techniques. Ultimately, \textsc{QLoRA} stands as a leveling mechanism, bridging the resource divide that exists between large industry players and smaller research groups utilizing consumer hardware.

An additional avenue for significant impact lies in the adaptation of LLM finetuning for mobile devices. Our \textsc{QLoRA} approach represents a potential breakthrough in enabling the personalization of LLMs directly on smartphones and other devices with constrained computational resources. While previous work has demonstrated the feasibility of running 7B parameter models on mobile hardware, \textsc{QLoRA} is the first solution that allows for efficient finetuning of such large models on these devices. Our projections suggest that an iPhone 12 Plus, for instance, could process up to 3 million tokens overnight during charging using \textsc{QLoRA}. Although the finetuned 7B models do not match the performance levels of ChatGPT, we argue that their capabilities are sufficiently robust to unlock novel use-cases previously hampered by either concerns over privacy or insufficient LLM performance. By leveraging \textsc{QLoRA}, users gain more control over their data and models, supporting privacy-focused applications and streamlining the deployment of LLMs.

Nonetheless, it is important to acknowledge that finetuning technologies possess dual-use characteristics and can be misappropriated for harmful purposes. There are documented risks associated with the widespread deployment of LLMs \citep{bommasani2021opportunities,bender2021dangers}, yet we contend that broadening public access to this rapidly proliferating technology can foster greater independent oversight and analysis, as opposed to restricting such capabilities to major corporate entities that withhold models and codebases from public scrutiny.

In conclusion, we anticipate that \textsc{QLoRA} will make a substantial positive contribution by democratizing access to high-quality LLM finetuning, thereby expanding its availability and ease of use.